% s-237-p237-m002; original PDF 237, 238
\Aphorism{s-237-p237-m002}{330}
Wie weit der Dichter höchstens den Redner mitnehmen kann, wenn ihm nicht endlich die Inspiration versagen soll, kann man an \Sperr{Gottfried Keller} beobachten, den seine Rhetorik immer erst im letzten Augenblicke verlässt, ehe sie unmöglich wird. Kein Dichter, Goethe nicht ausgenommen, hat ein ähnliches Genügen am Spiel und Gespinnst der Worte — freilich der kraftvollsten, fruchtbarsten und klingendsten Worte, wie sie wohl dem Meister Jacob Grimm vorschwebten, wenn er von der selbstzeugenden und gebärenden Macht des deutschen Wortes träumte: aber genug, es sind Worte, und immer ist durch Worte eine höhere künstlerische Wirkung gefährdet. In diesem letzten Moment besinnt sich Keller, dann verstummt die ewig abschweifende, umrankende, motivirende und übermotivirende Redseligkeit und es steht, nackt und leuchtend, ein Bildniss aus edelstem Marmor da, wie ihn wenige Hände zu formen bekamen. In der Kunst höchster verklärter Situationen, wo von den Menschenseelen ein stilles Glänzen ausgeht und der Brunnen der Schönheit wieder einmal klingt wie am Morgen des ersten Tages, in diesem Können des Zartesten und Gewaltigsten ist Keller ganz Meister, wenn er sich nur erst Phantasie und Hände dazu frei gemacht hat. Es giebt viele Dichter, die feinere Sinne und besseren Geschmack haben als er, und man braucht nur an seinen jüngeren Landsmann \Sperr{Conrad Ferdinand Meyer} zu denken, um in Keller etwas wie einen Bauern- und Volksdichter zu sehen, einen Nachfahren des vortrefflichen, künstlerisch aber ganz primitiven Jeremias Gotthelf; trotzdem ist, glaube ich, der Dichter des philiströsen Martin Salander und der zahlreichen politischen Gesinnungstüchtigkeiten der Ursprünglichere von Beiden. Aber Meyer ist der Geschliffenere, die bessere Gesellschaft, und so ist es in Ordnung: wenn für Keller immer der Werdende dankbar sein wird, mag es Meyer Dem, der fertig ist, recht machen — er, der selber fertig wurde, der spät seinen Erstling opferte, um früh vom Schreiben, wenn auch nicht vom Leben, Abschied zu nehmen und eine Reihe präciser gedrängter Novellen zu hinterlassen, in denen das \Sperr{Problem des grossen Menschen}, in seinen verschiedenen Spectraltypen, auf unvergängliche Weise abgehandelt wird, daneben einen Band Gedichte, deren Ernst, Süssigkeit und Concentration gegen alle andere Lyrik ungerecht stimmen. Ein deutsches Reich für zwei solche Schweizer!
\par

